%! Measuring Center and Spread — Unit 2, Lesson 2.2 — Student Packet Cover Sheet
% PILOT SHAPE, Unit 2 timing (5 / 45 / 5 / 5). THREE packet rows — Warm-Up, Guided Notes &
% Practice, Homework. The homework is the individual practice, assigned at the close and done
% outside class, and it is SCORED, so its cell is a \blank{}, never NA.
\documentclass[12pt]{article}
\usepackage{saar-article}
\usepackage{saar-boxes}
\usepackage{ltablex}
\keepXColumns
% Measured banner: the band grows to fit the title block, so a title long enough to wrap grows
% the band rather than running out of it (the stats course's \coverbanner; saar-article does not
% carry one, so it is defined here — every cover copies this block verbatim).
\newsavebox{\bannerbox}
\newlength{\bannerht}
\newlength{\coverbannerpad}
\setlength{\coverbannerpad}{0.16in}       % breathing room above and below the text
\newcommand{\coverbanner}[2]{%
  \sbox{\bannerbox}{%
    \begin{minipage}{\dimexpr\paperwidth-1.4in\relax}%
      \centering
      {\color{white}\textbf{\LARGE \CourseName}}\\[4pt]
      {\color{frostmid}\large\textbf{#1}}\\[6pt]
      {\color{white}\textbf{\large #2}}%
    \end{minipage}}%
  \setlength{\bannerht}{\dimexpr\ht\bannerbox+\dp\bannerbox+2\coverbannerpad\relax}%
  \noindent\begin{tikzpicture}[remember picture, overlay]
    \fill[cerulean] (current page.north west)
      rectangle ([yshift=-\bannerht]current page.north east);
    \node[anchor=north, inner sep=0pt]
      at ([yshift=-\coverbannerpad]current page.north) {\usebox{\bannerbox}};
  \end{tikzpicture}\par
  % The text body starts 0.6in below the page edge (geometry top=0.6in), so reserve only what
  % the band overruns that, plus a gap under the band.
  \vspace*{\dimexpr\bannerht-0.6in+0.16in\relax}%
}

\begin{document}

% Full-bleed cerulean banner, sized to the title block.
\coverbanner{Unit 2}{Lesson 2.2 \quad Measuring Center and Spread}

\vspace{0.06in}
% NAMESTRIP: this is the ONLY component that carries the name row.
\namedateperiod
\vspace{0.18in}

% GRADUAL RELEASE: the vocabulary is taught in the guided notes, so the targets name it
% plainly. The standard codes (PS.DS.2a, 2b, 2e) stay in the teacher-facing plan.
\begin{learningtargetbox}
\small
\begin{enumerate}[leftmargin=1.5em, itemsep=5pt]
  \item \ldots{} find the \textbf{mean} ($\bar{x}$), \textbf{median}, and \textbf{mode} of a
        data set and say what each one means in context.
  \item \ldots{} find the \textbf{range}, the \textbf{quartiles}, and the \textbf{interquartile
        range (IQR)}, give the \textbf{five-number summary}, and explain the IQR as the spread
        of the middle half of the data.
  \item \ldots{} compute the \textbf{variance} and the \textbf{standard deviation} from a table
        of \textbf{deviations}, explain why the deviations are squared and why we divide by
        $n - 1$, and read $s$ from a calculator's 1-Var Stats screen.
  \item \ldots{} explain that the \textbf{standard deviation} measures how far values typically
        sit from their mean --- not how big the values are.
\end{enumerate}
\end{learningtargetbox}

\vspace{0.14in}

\begin{tocbox}
\small
\renewcommand{\arraystretch}{1.5}
\begin{tabularx}{\linewidth}{@{} c l X r @{}}
\rowcolor{frost}
\textbf{\#} & \textbf{Component} & \textbf{Description} & \textbf{Score} \\
\midrule
1 & Warm-Up & Route 5's dot plot read and put in order, one average, and five numbers
             added, squared, and square-rooted & \blank{1.2cm} \\
2 & Guided Notes \& Practice & Mean, median, and mode; range, quartiles, and IQR; variance and
             standard deviation from a deviations table; four bus routes that separate spread
             from size --- then two towns' temperatures we work together & \blank{1.2cm} \\
% Row 3 is the lesson's GRADED practice, assigned at the close. Packet pages are the default; a
% Desmos activity may replace them on a given day, decided at Close & Assign. The row and its
% score cell never change.
3 & Homework & Oak Street Pizza's delivery times: center, the middle half, the standard
             deviation, and why slower deliveries are not more spread out ---
             \textbf{scored, due the first class after two study halls} & \blank{1.2cm} \\
\midrule
  & \multicolumn{2}{r}{\textbf{Total}} & \blank{1.2cm} \\
\end{tabularx}
\end{tocbox}

\vspace{0.14in}

% Keep in Mind carries the lesson's CONTENT — the rule, the test, the trap — never its process.
\begin{remindbox}
\small
The \textbf{mean} $\bar{x} = \sum x / n$ and the \textbf{median} measure center; the
\textbf{IQR} $= Q_3 - Q_1$ measures the spread of the middle half; the \textbf{standard
deviation} $s$ measures how far values typically sit from their mean. To find $s$: deviations,
square, add, divide by $n - 1$, square root --- and the deviations always add to $0$, which is
why we square them, not a sign of no spread. \textbf{Watch out:} a bigger standard deviation does
\emph{not} mean bigger values. Long commutes that huddle close together have a small $s$, and
adding the same amount to every value moves the mean but leaves $s$ unchanged.
\end{remindbox}

\end{document}
